\section{电压必须在两个点之间测量}

万用表的两支表笔接到两个节点，屏幕显示的是两点之间的电压差。把黑表笔接到
电路选定的 GND，红表笔接到信号线上，才有“这个引脚是 3.3 V”的说法。
GND 是共同参考，不是神秘的电流垃圾桶，也不保证所有位置永远完全同电位。

电压的单位是伏特，写作 V。数字电路不会要求输入恰好等于 0 V 或 3.3 V，
而是划出低电平和高电平范围。落在两个范围之间的电压不能可靠解释，软件也
无法通过读取更多次修复一个违反电气规格的输入。

\section{电流描述电荷移动的快慢}

电流的单位是安培，写作 A。对本书的第一遍主线，可以先把它理解为“每秒有
多少电荷经过一个截面”。电源向电路提供能量，导线和元件形成闭合路径；只有
从电源出发又能回到电源的路径，稳态电流才可能持续存在。

电阻限制电流。欧姆定律
\[
  V=IR
\]
连接电压 \(V\)、电流 \(I\) 和电阻 \(R\)。若 3.3 V 电源驱动一只压降约
1.9 V 的 LED，希望电流约为 4 mA，则串联电阻可先估算为
\[
  R=\frac{3.3-1.9}{0.004}\approx350\,\Omega.
\]
实际设计还要检查元件容差和额定功率，相关方法放在硬件深化附录。

\section{开关把连续电压变成两种可用状态}

机械开关可以接通或断开导电路径。晶体管也能受一个输入控制另一条电流路径，
而且速度远高于机械开关。大量晶体管组合后，可以让输出在接近电源和接近 GND
之间切换，于是上一章的 0 和 1 获得了物理载体。

发送端必须保证低电平足够低、高电平足够高；接收端必须给出自己能够识别的
范围。两边共同决定接口是否兼容。只看“都是 3.3 V 芯片”仍不够，方向、
允许电流、上电状态和时序也必须匹配。

\section{电源、时钟和复位承担不同职责}

芯片的电源引脚提供能量，GND 提供参考与回流路径，时钟引脚提供重复边沿，
复位引脚把内部状态带到规定起点。它们不能互相替代。电源已经存在时，时钟
可能仍未稳定；时钟已经跳变时，复位也可能仍保持有效。

CPU 上电后的简化顺序是：
\[
  \text{电源进入范围}
  \rightarrow\text{时钟稳定}
  \rightarrow\text{释放复位}
  \rightarrow\text{读取第一条指令}.
\]
每个箭头都需要时间。真实电路中的电容、稳压器、晶振和监控芯片怎样保证这些
时间，放到硬件深化附录再计算。

\section{先学会沿闭合路径检查故障}

若 LED 不亮，先检查电源、极性、电阻和回流，不要直接认定软件有错。若 CPU
没有输出，也应先问电源、时钟、复位和存放指令的器件是否都在工作。这个顺序
会一直延伸到操作系统调试：从最后一个确定发生的动作继续向后检查。

本章没有要求读者先掌握微积分、自然对数或开关电源。第一次主线只保留理解
数字信号和启动所需的电压、电流、开关、电源、时钟与复位；RC、KCL、MOSFET、
PCB 回流和 ESD 会在附录中逐项展开。
